\documentclass[aps,prl,twocolumn,superscriptaddress,showpacs]{revtex4-2}
\usepackage{amsmath,amssymb,amsfonts}
\usepackage{graphicx}
\usepackage{hyperref}
\usepackage{xcolor}

\begin{document}

\title{Letter 71: Baryogenesis from BCT Vacuum Geometry\\
\large The Baryon Asymmetry as a Three-Loop Josephson Effect}

\author{Michel Robert Cabri\'{e}}
\affiliation{Independent Researcher, Victoria, Australia}
\email{ZeroFreeParameters@gmail.com}

\begin{abstract}
We derive the cosmic baryon asymmetry $\eta_B = n_B/s$ from the geometry 
of the BCT Superfluid Lattice vacuum, with zero free parameters. 
The OHC phase transition at $T_\mathrm{BCT} \approx 0.04\,T_P$ creates 
a net Hopf charge asymmetry through CP violation encoded in the D4 triality 
Jarlskog invariant $J_\mathrm{BCT} = \sqrt{3}/2$ (exact). This asymmetry 
is topologically protected through cosmic cooling and converted to baryon 
number at the electroweak transition via three successive Josephson phase 
slips --- one per fermion generation. The result is
$\eta_B = (45\sqrt{3}\times 28)/(4\pi^2 g_* \times 79)\times\alpha_0\times(\alpha\pi)^3
= 5.853\times 10^{-10}$,
within $-4.0\%$ of the observed value $6.10\times 10^{-10}$.
All inputs are BCT geometry. No new physics is required.
\end{abstract}

\maketitle

\section{Introduction}

The observed baryon asymmetry of the universe,
\begin{equation}
\eta_B \equiv \frac{n_B}{s} = (6.10 \pm 0.04)\times 10^{-10},
\end{equation}
remains one of the deepest unsolved problems in physics. The Standard Model 
satisfies the three Sakharov conditions~\cite{Sakharov1967} only marginally, 
and cannot quantitatively account for the observed asymmetry. Every proposed 
solution requires physics beyond the Standard Model.

We show here that the BCT Superfluid Lattice Model~\cite{BCT_Monograph} 
provides a complete, parameter-free derivation of $\eta_B$ from vacuum 
geometry. The mechanism requires no new particles, no extensions to the 
Standard Model, and no free parameters beyond the three geometric inputs 
of the BCT programme: $r_\mathrm{oct}$, $r_\mathrm{tet}$, and 
$\Lambda_\mathrm{QCD} = 220\,\mathrm{MeV}$.

\section{The BCT Vacuum and the OHC Ground State}

The BCT vacuum ground state is the Octet-Hopfion Condensate (OHC)~\cite{BCT_AppJH},
a topological superfluid condensate filling each BCT Planck sphere of radius 
$R_s = 12\,l_P$ with Hopf charge $H \in \mathbb{Z}$. The ground state carries 
$H = 1$, topologically protected by $\pi_3(S^2) = \mathbb{Z}$.

The two canonical BCT void radii at $c/a = \sqrt{2}$ are:
\begin{align}
r_\mathrm{oct} &= \frac{\sqrt{2}-1}{2} \approx 0.207107, \\
r_\mathrm{tet} &= \frac{\sqrt{6}-2}{4} \approx 0.112372.
\end{align}
These define the bare coupling and fine structure constant:
\begin{align}
\alpha_0 &= \frac{r_\mathrm{oct}\times r_\mathrm{tet}}{\pi} = 0.0074081, \\
\alpha &= \alpha_0(1-2\alpha_0) \Rightarrow 1/\alpha = 137.018.
\end{align}

\section{The BCT Cosmic Timeline}

The BCT vacuum undergoes a hierarchy of phase transitions:

\begin{enumerate}
\item \textbf{$T = T_\mathrm{BCT} \approx 0.04\,T_P$}: OHC condensation. 
The interior condensate $\Psi_2$ acquires $H = +1$ winding. D4 triality 
CP asymmetry biases $H = +1$ over $H = -1$.

\item \textbf{$T_\mathrm{BCT} \to T_\mathrm{EW}$}: Topological freeze. 
The net Hopf charge asymmetry is \emph{exactly preserved} by topology --- 
$H \in \mathbb{Z}$ cannot change under smooth evolution of $\Psi_2$.

\item \textbf{$T = T_\mathrm{EW} = 246.145\,\mathrm{GeV}$}: Electroweak 
symmetry breaking. Weak sphalerons convert Hopf charge to baryon number 
via three Josephson phase slips.

\item \textbf{$T = \Lambda_\mathrm{QCD} = 220\,\mathrm{MeV}$}: QCD 
confinement. Baryon asymmetry locked into hadrons.
\end{enumerate}

\section{The Three Sakharov Conditions in BCT}

All three Sakharov conditions~\cite{Sakharov1967} are satisfied by BCT geometry 
without additional assumptions:

\textbf{(1) Baryon number violation.} BCT sphaleron transitions at the Planck 
scale violate baryon number through topological Hopf charge flips. The rate is 
set by $\exp(-S_\mathrm{D4})$ where $S_\mathrm{D4} = \pi^5/6 = 51.003$.

\textbf{(2) C and CP violation.} The D4 lattice triality is not symmetric under 
CP. The three OHC triality states carry phases $\phi_k = e^{2\pi i k/3}$, $k=0,1,2$. 
The Jarlskog-like CP invariant is:
\begin{equation}
J_\mathrm{BCT} = \mathrm{Im}[\phi_0\,\phi_1^*\,\phi_2] = \frac{\sqrt{3}}{2}
\quad\text{(exact).}
\end{equation}
The CP phase is maximal: $\delta_\mathrm{CP} = -\pi/2$, consistent with the 
BCT leptogenesis prediction of Letter~53~\cite{BCT_L53}.

\textbf{(3) Departure from thermal equilibrium.} The OHC condensation at 
$T_\mathrm{BCT}$ is a first-order phase transition. The $H=0 \to H=1$ 
condensation is sudden, departing from thermal equilibrium.

\section{The CP Asymmetry and Topological Protection}

The net CP asymmetry per BCT sphere at condensation is:
\begin{equation}
\epsilon_\mathrm{CP} = \alpha_0 \times J_\mathrm{BCT} 
= \alpha_0 \times \frac{\sqrt{3}}{2} = 6.416\times 10^{-3}.
\end{equation}

\textbf{Topological Protection Theorem.} The Hopf invariant $H \in \mathbb{Z}$ 
is a homotopy invariant of $\pi_3(S^2) = \mathbb{Z}$. It cannot change under 
any smooth deformation of $\Psi_2$ that is continuous away from the sphere 
boundary. Changing $H$ requires passing through a nodal surface where 
$|\Psi_2| = 0$, which costs infinite energy in the superfluid. 
\emph{Therefore, the net Hopf charge asymmetry created at $T_\mathrm{BCT}$ 
is exactly preserved through all subsequent cosmic evolution until $T_\mathrm{EW}$.}

The number density of Hopf charge asymmetry per entropy unit is:
\begin{equation}
\frac{n_H}{s} = \epsilon_\mathrm{CP} \times \frac{45}{2\pi^2 g_*},
\end{equation}
where $g_* = 106.75$ counts relativistic degrees of freedom at $T_\mathrm{EW}$.

\section{Three-Loop Josephson Conversion at the Electroweak Scale}

At $T_\mathrm{EW}$, the electroweak sphalerons convert Hopf charge to 
baryon number. In the BCT framework, this conversion proceeds through 
\emph{three successive Josephson phase slips} across the OHC boundary 
--- one per fermion generation, one per D4 triality state.

Each Josephson phase slip carries amplitude $\sim \alpha\pi$: the factor 
$\alpha$ from the electromagnetic coupling strength at the sphere surface, 
and $\pi$ from the geometric phase accumulated traversing the Hopf fibre.  
Three slips give the suppression factor:
\begin{equation}
f_J = (\alpha\pi)^3 = (7.298\times 10^{-3} \times \pi)^3 = 1.205\times 10^{-5}.
\end{equation}

The electroweak sphaleron conversion efficiency is given by the 
Khlebnikov--Shaposhnikov result~\cite{KhlebShap1988}:
\begin{equation}
f_\mathrm{sph} = \frac{28}{79} = 0.3544.
\end{equation}

\section{The Baryon Asymmetry}

Combining all factors:
\begin{equation}
\boxed{
\eta_B = \frac{45\sqrt{3}\times 28}{4\pi^2\, g_*\times 79}
\times \alpha_0 \times (\alpha\pi)^3
}
\end{equation}

Substituting BCT values:
\begin{align}
\eta_B &= \frac{45\times\frac{\sqrt{3}}{2}\times 28 \times 2}
{4\pi^2 \times 106.75 \times 79}
\times 0.0074081 \times (0.007298\times\pi)^3 \nonumber\\
&= 5.853\times 10^{-10}.
\end{align}

\begin{table}[h]
\centering
\caption{BCT baryogenesis prediction vs observation}
\begin{tabular}{lll}
\hline\hline
Quantity & Value & Source \\
\hline
$\eta_B$ (BCT) & $5.853\times 10^{-10}$ & This work \\
$\eta_B$ (Planck 2018) & $6.100\times 10^{-10}$ & \cite{Planck2018} \\
Error & $-4.0\%$ & \\
\hline\hline
\end{tabular}
\end{table}

\section{All Inputs are BCT Geometry}

Every factor in Eq.~(9) is determined by BCT void geometry with zero free parameters:
\begin{align}
\alpha_0 &= r_\mathrm{oct}\times r_\mathrm{tet}/\pi 
\quad\text{[bare coupling]},\\
\alpha &= \alpha_0(1-2\alpha_0) 
\quad\text{[renormalised coupling]},\\
\sqrt{3}/2 &= J_\mathrm{BCT} 
\quad\text{[D4 triality, exact]},\\
28/79 &\quad\text{[EW sphaleron conversion, standard result]}.
\end{align}

The quantity $g_* = 106.75$ counts the Standard Model degrees of freedom, 
themselves derived from the BCT gauge group emergence~\cite{BCT_L18}.

\section{Discussion}

The BCT derivation of $\eta_B$ completes the programme's account of cosmic 
evolution. The same geometric quantities that determine the fine structure 
constant ($\alpha_0$, $\alpha$) and the three fermion generations (D4 triality, 
$J_\mathrm{BCT} = \sqrt{3}/2$) also determine why matter dominates antimatter 
in the observed universe.

The $-4.0\%$ discrepancy is consistent with the known theoretical uncertainty 
in the Khlebnikov--Shaposhnikov sphaleron conversion factor and the effective 
$g_*$ at the electroweak scale. A complete treatment of BCT degrees of freedom 
at $T_\mathrm{BCT}$, deferred to Appendix JH5, may improve this precision.

The central result stands: the baryon asymmetry of the universe is a 
three-loop Josephson effect in the BCT vacuum geometry. You exist because 
$(\alpha\pi)^3 = 1.205\times 10^{-5}$.

\begin{acknowledgments}
The author thanks the Zenodo open-access platform. 
This result was obtained in Melbourne, Australia, 17 March 2026,
over the last of the birdseed money.
\end{acknowledgments}

\begin{thebibliography}{99}

\bibitem{Sakharov1967}
A.~D. Sakharov,
\textit{JETP Lett.} \textbf{5}, 24 (1967).

\bibitem{BCT_Monograph}
M.~R. Cabri\'{e},
``The BCT Superfluid Lattice Model,''
\href{https://doi.org/10.5281/zenodo.18884976}{doi:10.5281/zenodo.18884976} (2026).

\bibitem{BCT_AppJH}
M.~R. Cabri\'{e},
``BCT Appendix JH: The Hopfion Interior and the Octet-Hopfion Condensate,''
\href{https://doi.org/10.5281/zenodo.18975018}{doi:10.5281/zenodo.18975018} (2026).

\bibitem{BCT_L18}
M.~R. Cabri\'{e},
``BCT Letter 18: The D4 Uniqueness Theorem,''
\href{https://doi.org/10.5281/zenodo.18957202}{doi:10.5281/zenodo.18957202} (2026).

\bibitem{BCT_L53}
M.~R. Cabri\'{e},
``BCT Letter 53: Maximal CP Violation and Leptogenesis,''
BCT Programme (2026).

\bibitem{KhlebShap1988}
S.~Yu. Khlebnikov and M.~E. Shaposhnikov,
\textit{Nucl. Phys.} \textbf{B308}, 885 (1988).

\bibitem{Planck2018}
Planck Collaboration,
\textit{Astron. Astrophys.} \textbf{641}, A6 (2020).

\end{thebibliography}

\end{document}
